% TC5 / Group B smoke-test artifact for the pre-review-notes skill.
\documentclass[a4paper,10pt]{article}

\usepackage[utf8]{inputenc}
\usepackage[left=2.2cm,right=1.2cm,top=1.2cm,bottom=1.2cm]{geometry}
\usepackage{xcolor}
\usepackage{amsmath,amssymb}
\usepackage{tikz}
\usepackage{pifont}
\usepackage{tcolorbox}
\usepackage{hyperref}
\usepackage{booktabs}
\usepackage{everypage}

\tcbuselibrary{skins,breakable}

\definecolor{barblue}{HTML}{2E75B6}
\definecolor{titleblue}{HTML}{2E75B6}
\definecolor{redline}{HTML}{C0392B}
\definecolor{yellowbg}{HTML}{FFF3CD}
\definecolor{yellowtext}{HTML}{856404}
\definecolor{lightblue}{HTML}{D6EAF8}

\pagestyle{empty}
\setlength{\parindent}{0pt}
\setlength{\emergencystretch}{2em}
\hypersetup{colorlinks=true,linkcolor=titleblue,urlcolor=titleblue}

\newcommand{\chaptertitle}[1]{{\large\bfseries\color{titleblue}#1}\\[2pt]}
\newcommand{\pref}[1]{\colorbox{yellowbg}{\scriptsize\texttt{\color{yellowtext}#1}}}
\newcommand{\subtitle}[1]{{\scriptsize\color{gray}\itshape#1}}
\newcommand{\topicsection}[1]{%
  \vspace{6pt}{\noindent\normalsize\bfseries\color{titleblue}#1}\\[-6pt]
  {\color{barblue}\rule{\linewidth}{1.2pt}}\\[2pt]
}

\newtcolorbox{rc}{
  blanker,breakable,
  borderline west={2pt}{0pt}{redline},
  left=8pt,top=1pt,bottom=1pt,
  before skip=0pt,after skip=3pt,
}
\newtcolorbox{keybox}{
  colback=lightblue,colframe=barblue,
  boxrule=0.8pt,arc=2pt,
  left=4pt,right=4pt,top=2pt,bottom=2pt,
  before skip=4pt,after skip=4pt,
}
\newcommand{\step}[2]{{\small\color{redline}\bfseries#1}\;{\small#2}\\[1pt]}

\AddEverypageHook{\tikz[remember picture,overlay]{
  \fill[barblue] ([yshift=0pt]current page.north west)
    rectangle ([yshift=-14pt]current page.north east);
}}

\begin{document}

\vspace*{2pt}
\chaptertitle{Pilot Chapter: Regression}
\pref{02-regression.pdf pp. 2--34}\quad
\subtitle{Selected as a prerequisite-level smoke test, not as an assessment-priority claim.}

\vspace{4pt}

\begin{keybox}
\small\textbf{Evidence boundary.} The supplied material contains lecture slides only.
No grading breakdown, exam format, past paper, marking rubric, or official solution was
available. Accordingly, this chapter makes no claim that regression is highly weighted,
frequently tested, or representative of a current STAT3007 offering.
\end{keybox}

\topicsection{1. Learning setup}
\pref{02-regression.pdf pp. 2, 7, 14}

\begin{rc}
In supervised learning, the data are pairs $(x_i,y_i)$. A regression target $y_i$
is real-valued, whereas a classification target is discrete. The lecture frames model
learning as four linked choices:
\begin{enumerate}
  \item collect training data;
  \item choose a hypothesis or model class;
  \item choose an objective that measures fit;
  \item optimize the objective to obtain model parameters.
\end{enumerate}
Training fit is not the final goal: the learned model should also work on unseen data.
Regularization injects prior structure or discourages fitting irregularities in the sample.

\begin{keybox}
\centering\small\bfseries
A regression method is defined jointly by its model class, objective, optimization
procedure, and generalization assumptions.
\end{keybox}
\end{rc}

\topicsection{2. Ordinary least squares}
\pref{02-regression.pdf pp. 16--19}

\begin{rc}
For a design matrix $X\in\mathbb{R}^{n\times d}$ with row $x_i^\top$ and response
vector $y\in\mathbb{R}^n$, ordinary least squares (OLS) minimizes
\[
  R_n(\beta)=\sum_{i=1}^{n}(x_i^\top\beta-y_i)^2
             =\lVert X\beta-y\rVert_2^2.
\]
The normal equations are
\[
  X^\top X\beta=X^\top y.
\]
If $X^\top X$ is invertible, the unique solution is
\[
  \widehat\beta_{\mathrm{OLS}}=(X^\top X)^{-1}X^\top y.
\]
If $X^\top X$ is singular, the inverse formula is invalid and the normal equations can
have multiple solutions. Under Gaussian conditional noise with fixed variance, maximizing
the conditional likelihood is equivalent to minimizing the same squared-error objective.
\end{rc}

\topicsection{Worked derivation: the OLS normal equations}
\pref{02-regression.pdf p. 18}

\begin{rc}
\textbf{Source problem.} Derive the closed-form OLS estimator by differentiating the
sum of squared errors.

\medskip
\textbf{Independently checked solution:}

\step{\ding{192}}{Expand the quadratic:
$R_n(\beta)=\beta^\top X^\top X\beta-2y^\top X\beta+y^\top y$.}
\step{\ding{193}}{Differentiate. Since $X^\top X$ is symmetric,
$\nabla_\beta R_n(\beta)=2X^\top X\beta-2X^\top y$.}
\step{\ding{194}}{Set the gradient to zero to obtain
$X^\top X\beta=X^\top y$.}
\step{\ding{195}}{When $X^\top X$ is invertible, left-multiply by its inverse:
$\widehat\beta=(X^\top X)^{-1}X^\top y$.}

\begin{keybox}
\centering\small\bfseries
The inverse formula needs a full-column-rank design; the normal equations themselves do not.
\end{keybox}
\end{rc}

\topicsection{3. Ridge regression and the bias term}
\pref{02-regression.pdf pp. 20--25}

\begin{rc}
Ridge regression adds an $\ell_2$ penalty:
\[
  \widehat\beta_{\mathrm{ridge}}
  =\arg\min_{\beta}\left(\lVert X\beta-y\rVert_2^2
  +\lambda\lVert\beta\rVert_2^2\right),\qquad \lambda>0,
\]
which gives
\[
  \widehat\beta_{\mathrm{ridge}}
  =(X^\top X+\lambda I)^{-1}X^\top y.
\]
The solution is unique because, for every nonzero $v$,
\[
  v^\top(X^\top X+\lambda I)v
  =\lVert Xv\rVert_2^2+\lambda\lVert v\rVert_2^2>0.
\]
Thus $X^\top X+\lambda I$ is positive definite. With an intercept $b$, the lecture
penalizes the slope coefficients but not $b$:
\[
  \min_{b,\beta}\sum_i(x_i^\top\beta+b-y_i)^2
  +\lambda\lVert\beta\rVert_2^2.
\]
Centering $x_i$ and $y_i$ reduces this to ridge regression without an intercept, after
which $\widehat b=\bar y-\bar x^\top\widehat\beta$.
\end{rc}

\topicsection{4. Nonlinear and nonparametric regression}
\pref{02-regression.pdf pp. 26--32}

\begin{rc}
\textbf{Basis functions.} Map $x$ to
$(g_1(x),\ldots,g_k(x))$ and fit a linear model in the transformed features. The model
is linear in its learned coefficients even when it is nonlinear in the original input.

\medskip
\textbf{$k$-nearest-neighbour regression.} For a query $x$,
\[
  h_n(x)=\frac{1}{k}\sum_{x_i\in N_k(x)}y_i.
\]
Its effective representation grows with the stored data, so it is nonparametric in the
course's sense; this does not mean it has no tunable choices.

\medskip
\textbf{Kernel regression.} With a similarity function $K$,
\[
  h_n(x)=\frac{\sum_i K(x,x_i)y_i}{\sum_i K(x,x_i)}.
\]
Both nearest-neighbour and kernel methods can represent flexible functions, but the
lecture notes their increasing data and computational demands, including the curse of
dimensionality for nearest neighbours.
\end{rc}

\topicsection{Worked exercise: the 1-NN prediction function}
\pref{02-regression.pdf pp. 27--29}

\begin{rc}
\textbf{Problem statement (from the slide):} ``Given three training examples with
$(x,y)$ being $(1,1)$, $(2,2)$, $(4,4)$, draw the function learned by 1NN.''

\medskip
\textbf{Independently checked solution:}

\step{\ding{192}}{The nearest-neighbour boundaries are the midpoints $1.5$ and $3$.}
\step{\ding{193}}{Away from ties,
$h(x)=1$ for $x<1.5$, $h(x)=2$ for $1.5<x<3$, and $h(x)=4$ for $x>3$.}
\step{\ding{194}}{At $x=1.5$ and $x=3$, two samples are equally near; the prediction
depends on the implementation's tie rule.}
\step{\ding{195}}{Adding $(5,5)$ introduces a new boundary at $4.5$ and a final region
with prediction $5$. Adding $(3,3)$ instead gives boundaries $1.5$, $2.5$, and $3.5$.}

\begin{keybox}
\centering\small\bfseries
A 1-NN regressor is piecewise constant, with transitions at Voronoi boundaries.
\end{keybox}
\end{rc}

\topicsection{\color{redline}\ding{72} Verification finding: a source inconsistency}
\pref{02-regression.pdf pp. 27--28}

\begin{rc}
Page 27 defines 1-NN prediction as the response of the nearest training point. That
definition yields a piecewise-constant function. Page 28 describes the displayed 1-NN
fit as piecewise-linear. These statements are inconsistent unless an unstated interpolation
rule is being used. The worked exercise above follows the explicit definition on page 27
and flags ties rather than silently copying the page-28 wording.

\begin{keybox}
\centering\small\bfseries\color{redline}
When prose conflicts with a displayed definition, recompute from the definition and record
the discrepancy.
\end{keybox}
\end{rc}

\topicsection{Worked exercise: end-of-lecture checks}
\pref{02-regression.pdf p. 33}

\begin{rc}
The slide asks which statements are correct. The independently checked answers are:
\begin{itemize}
  \item Penalizing a large bias term: \textbf{false} for the ridge formulation used here;
        page 25 leaves the intercept unpenalized.
  \item Ridge always has a unique solution: \textbf{true} under the stated $\lambda>0$;
        the positive-definiteness argument above proves it.
  \item A nonparametric method has no parameters: \textbf{false}; ``nonparametric'' here
        means model capacity can grow with the data.
  \item Learning is often formulated as optimization: \textbf{true}; this is the
        course's framing on pages 7 and 14.
\end{itemize}
Therefore the correct choices are the second and fourth statements.
\end{rc}

\topicsection{5. Quick reference card}
\pref{02-regression.pdf p. 34 and cited pages above}

\begin{rc}
\textbf{OLS objective:} $\lVert X\beta-y\rVert_2^2$.

\textbf{OLS solution:} $(X^\top X)^{-1}X^\top y$ only when $X^\top X$ is invertible.

\textbf{Ridge objective:} OLS $+\lambda\lVert\beta\rVert_2^2$, $\lambda>0$.

\textbf{Ridge solution:} $(X^\top X+\lambda I)^{-1}X^\top y$; unique for $\lambda>0$.

\textbf{Intercept:} add a dummy feature for OLS; for the course's ridge formulation,
center the data and do not penalize the intercept.

\textbf{Basis method:} transform inputs, then fit a model linear in the new coefficients.

\textbf{$k$NN:} average the responses of the $k$ nearest points.

\textbf{Kernel regression:} similarity-weighted average of all observed responses.

\begin{keybox}
\centering\small\bfseries\color{redline}
Check rank assumptions, distinguish an intercept from penalized slopes, and remember that
1-NN predictions are piecewise constant under the stated definition.
\end{keybox}
\end{rc}

\topicsection{6. Evidence and practice mapping}

\begin{rc}
\small
\textbf{Assessment evidence:} none was present in the downloaded slide set. The table
therefore maps instruction and in-lecture practice only; it is not an exam map.

\medskip
\renewcommand{\arraystretch}{1.25}
\begin{tabular}{@{}lp{10.2cm}@{}}
\toprule
\textbf{Source} & \textbf{Role in this pilot chapter} \\
\midrule
02-regression.pdf p. 18 & Instructor derivation of the OLS closed form; independently re-derived. \\
02-regression.pdf p. 29 & Instructor 1-NN drawing exercise; solved with tie cases stated. \\
02-regression.pdf p. 33 & Instructor multiple-choice check; each option justified. \\
02-regression.pdf p. 34 & Explicit ``What You Need to Know'' content boundary for the lecture. \\
\bottomrule
\end{tabular}

\begin{keybox}
\centering\small\bfseries
Before calling any item ``exam-important'', add the current course profile, assessment
weighting, past papers, and official solutions.
\end{keybox}
\end{rc}

\end{document}
